%=============================================================================
% WAVE 4, ITERATION 16: Q_psi LOCAL vs GLOBAL --- THE DEFINITIVE ANALYSIS
%
% Agent: P2 Specialist (Wave 4 i16, Opus 4.6) | DFD Research Programme
% Date: 20 March 2026
%
% Building on: W4i15_P2_update.tex (omega^2 derivation, Q_psi ~30 estimate),
%              W4i15_P11_update.tex (overlap-integral derivation, C7 shape),
%              section02_formalism.tex (action principle, field equations),
%              appendix_Q_temporal_completion.tex (dust branch, c_s^2 -> 0),
%              appendix_R_em_psi_coupling.tex (mode equation, Omega_psi, gamma_psi),
%              Deep_EM_Coupling_Analysis.tex (modal decomposition, channels 1+2),
%              MASTER_FINDINGS_CORPUS.md (Q_psi = 10^9, MOND self-confinement)
%
% THE QUESTION: Is the relevant Q_psi for SRF cavity loss the LOCAL cavity
%   psi-mode Q (~30, radiation-limited) or the GLOBAL MOND self-confinement
%   value (10^9)? If local, the SQMS bound tightens by ~5000x.
%
% Status flags: [VERIFIED], [NEW], [CORRECTED], [CAUTION], [HONEST]
%=============================================================================

\documentclass[11pt]{article}
\usepackage[utf8]{inputenc}
\usepackage{amsmath,amssymb,amsfonts,amsthm,mathtools}
\usepackage{geometry}
\geometry{margin=1in}
\usepackage{booktabs}
\usepackage{siunitx}
\usepackage{hyperref}
\usepackage{xcolor}
\usepackage{enumitem}
\usepackage{float}
\usepackage{array}
\usepackage{tcolorbox}
\tcbuselibrary{skins,breakable}

\newcommand{\dpsi}{\delta\psi}
\newcommand{\lam}{\lambda}
\newcommand{\lbone}{|\lambda - 1|}
\newcommand{\Qpsi}{Q_\psi}
\newcommand{\Meff}{M_{\rm eff}}
\newcommand{\ee}[1]{\times 10^{#1}}
\newcommand{\half}{\tfrac{1}{2}}
\newcommand{\kB}{k_{\mathrm{B}}}
\newcommand{\Opsi}{\Omega_\psi}
\newcommand{\gpsi}{\gamma_\psi}
\newcommand{\Mpsi}{M_\psi}
\newcommand{\order}[1]{\mathcal{O}(#1)}
\newcommand{\psigrav}{\psi_{\rm grav}}
\newcommand{\dpsiEM}{\delta\psi_{\rm EM}}
\newcommand{\calF}{\mathcal{F}}
\newcommand{\calB}{\mathcal{B}}
\newcommand{\Heff}{H_n}

\definecolor{strongcolor}{rgb}{0,0.45,0}
\definecolor{weakcolor}{rgb}{0.65,0,0}
\definecolor{conjcolor}{rgb}{0.6,0.3,0}

\newcommand{\confirmed}[1]{\textcolor{blue}{\textbf{CONFIRMED:} #1}}
\newcommand{\corrected}[1]{\textcolor{red}{\textbf{CORRECTED:} #1}}
\newcommand{\caution}[1]{\textcolor{orange}{\textbf{CAUTION:} #1}}
\newcommand{\honest}[1]{\textcolor{violet}{\textbf{HONEST ASSESSMENT:} #1}}
\newcommand{\verdict}[1]{\textcolor{green!50!black}{\textbf{VERDICT:} #1}}
\newcommand{\newfinding}[1]{\textcolor{teal}{\textbf{NEW FINDING:} #1}}

\newtheorem{theorem}{Theorem}[section]
\newtheorem{proposition}[theorem]{Proposition}
\newtheorem{corollary}[theorem]{Corollary}
\newtheorem{lemma}[theorem]{Lemma}
\theoremstyle{definition}
\newtheorem{definition}[theorem]{Definition}
\newtheorem{finding}[theorem]{Finding}
\newtheorem{correction}[theorem]{Correction}
\theoremstyle{remark}
\newtheorem{remark}[theorem]{Remark}

\title{Wave 4, Iteration 16: $Q_\psi$ Local vs.\ Global\\
\large The Definitive Analysis of Scalar-Mode Quality Factor\\
in SRF Cavity Loss and Its Impact on SQMS Phase~I}

\author{DFD Research Programme --- P2 Agent (W4i16, Opus 4.6)\\
\textit{Building on: W4i15 P2/P11 derivations, section02\_formalism,\\
appendix\_Q (temporal completion), appendix\_R (EM-$\psi$ coupling)}}

\date{20 March 2026}

\begin{document}
\maketitle
\tableofcontents
\newpage

%=============================================================================
\section*{Executive Summary}
%=============================================================================

\begin{tcolorbox}[colback=blue!5!white, colframe=blue!60!black,
  title=\textbf{W4i16 --- $Q_\psi$ LOCAL vs GLOBAL: RESOLUTION}]

The W4i15 P2 derivation of the $\omega^2$ psi-channel loss scaling
raised a critical question: is the relevant $\Qpsi$ in the loss formula
the \textbf{local} cavity psi-mode quality factor ($\Qpsi^{\rm local}
\sim 30$, set by radiation losses from the cavity volume) or the
\textbf{global} MOND self-confinement value ($\Qpsi^{\rm MOND} \sim
10^9$)? The difference is a factor $\sim 3\ee{7}$, which would change
the psi-channel loss rate by the same factor and tighten the SQMS bound
by $\sim 5000\times$.

\textbf{The answer is: NEITHER in the simple form stated.} The correct
analysis reveals a three-regime structure that depends on the $c_s^2
\to 0$ dust branch theorem and the two-component decomposition $\psi
= \psigrav + \dpsiEM$.

\begin{enumerate}[nosep]
\item \textbf{FINDING 1 [NEW, PROGRAMME-CRITICAL]:} The dust branch
  ($c_s^2 \to 0$) from Appendix~Q of the main text means the
  \emph{free} psi-field has zero propagation speed in the cosmological
  limit. A mode with $c_s = 0$ \textbf{cannot radiate}: radiation
  requires finite group velocity to carry energy away from the source
  region. Therefore the W4i15 estimate $\Qpsi^{\rm local} \sim
  (k R)^3 \sim 30$ is \textbf{WRONG} --- it assumed $c_\psi = c$
  (the Newtonian-limit propagation speed), but the correct DFD
  propagation speed is $c_s \to 0$ on the dust branch.

\item \textbf{FINDING 2 [NEW, PROGRAMME-CRITICAL]:} However, $c_s^2
  \to 0$ applies to the \emph{cosmological} temporal sector. In the
  \emph{spatial} sector (laboratory, quasi-static), the field equation
  is elliptic (Poisson-like) with $\mu \to 1$. The ``psi-mode'' inside
  the cavity is not a propagating wave --- it is an instantaneous
  Coulomb-like response to the EM source. The correct object is not
  a resonant mode with a $Q$-factor, but a \textbf{driven, non-resonant,
  quasi-static response}. The Green's function is $\hat{G} \sim 1/k^2$,
  not a Lorentzian resonance.

\item \textbf{FINDING 3 [CORRECTED, PROGRAMME-CRITICAL]:} Re-deriving
  the loss from the quasi-static (non-resonant) response eliminates
  $\Qpsi$ from the formula entirely. The psi-channel loss becomes:
  \begin{equation}
  \frac{1}{Q_\psi^{\rm cav}(\omega)} =
  \frac{(\lambda - 1)^2\,\Heff^2(\omega)}
  {\omega}\,\frac{4\pi G\,U_{\rm EM}}{c^4\,\Mpsi\,\Opsi^2},
  \tag{$\star$}
  \end{equation}
  where $\Mpsi\,\Opsi^2 = a_*^2/(8\pi G) \times k^2 V$ is the
  static stiffness of the psi-field mode, with \textbf{no $\Qpsi$
  denominator}. The Q-ratio prediction is unchanged: $\mathcal{R}
  = 0.49$ (overlap-integral ratio) or $0.27$ (with $1/\omega$
  correction). The SQMS bound $\lbone < 1.56\ee{-9}$ requires
  reinterpretation.

\item \textbf{FINDING 4 [NEW, HIGH IMPACT]:} The reinterpretation:
  in all prior iterations, the SQMS bound was derived assuming
  $\Qpsi = 10^9$ appeared in the denominator of the loss formula.
  If $\Qpsi$ drops out (quasi-static response), the predicted
  loss at given $\lbone$ is \textbf{larger} by a factor $\sim
  \Qpsi^{\rm MOND}/1 = 10^9$. Equivalently, the SQMS bound
  \textbf{tightens} by $\sqrt{10^9} \approx 3.2\ee{4}$:
  \begin{equation}
  \lbone < \frac{1.56\ee{-9}}{\sqrt{10^9}} \approx 5\ee{-14}.
  \tag{$\star\star$}
  \end{equation}
  \caution{This is a factor $\sim 30{,}000\times$ tighter than the
  previous bound. If correct, it means existing SQMS data ALREADY
  constrain $\lbone$ to the $10^{-14}$ level, making the Phase~I
  ``intentional'' search largely redundant --- the accidental bound
  from cavity stability would already be at the target sensitivity.}

\item \textbf{FINDING 5 [HONEST ASSESSMENT]:} The resolution of this
  tension requires understanding WHICH regime applies: the fully
  cosmological ($c_s^2 = 0$) regime, the Newtonian ($\mu = 1$,
  $c_\psi = c$) regime, or a transitional regime. The laboratory
  cavity sits deep in the Newtonian regime for the \emph{spatial}
  sector, but the \emph{temporal} response depends on $K(\Delta)$
  which gives $c_s^2 \to 0$. This creates an asymmetry: spatial
  Laplacian with $\mu = 1$ (fast, Coulomb-like) vs.\ temporal
  derivative with $c_s^2 \to 0$ (slow, dust-like). The psi-mode
  in the cavity is then \textbf{spatially Newtonian but temporally
  frozen} --- a quasi-static mode with infinite effective $Q$.
\end{enumerate}

\verdict{The local vs.\ global $\Qpsi$ question is resolved by
recognizing that neither applies. The correct regime is quasi-static
(no resonant $Q$), set by the $c_s^2 \to 0$ dust branch combined with
Newtonian-limit spatial response. This eliminates $\Qpsi$ from the
loss formula and potentially tightens the SQMS bound by $\sim
3\ee{4}\times$. \textbf{This is programme-changing and requires
immediate independent verification.}}
\end{tcolorbox}


%=============================================================================
%=============================================================================
\part{The Three Candidate Regimes}
\label{part:regimes}
%=============================================================================
%=============================================================================


%=============================================================================
\section{Regime 1: Global MOND Self-Confinement ($\Qpsi = 10^9$)}
\label{sec:regime1}
%=============================================================================

The $\Qpsi = 10^9$ value appears throughout the programme (MASTER\_FINDINGS\_CORPUS
line 717: ``$Q_\psi$ derived: $10^9$, MOND self-confinement''; W4i15 P2 inherited
fact line 15: ``$Q_\psi = 10^9$ (MOND self-confinement)'').

\subsection{Origin of the $10^9$ estimate}

The global $\Qpsi$ arises from treating the galactic-scale psi-field
as a self-gravitating mode confined by the MOND nonlinearity. In the
deep-field regime ($\mu \sim x$), the nonlinear field equation
supports quasi-bound ``solitonic'' solutions where the field amplitude
is maintained by the self-interaction term $\kappa a^2/c^2$. The
quality factor of such modes is set by the ratio of stored energy to
radiation loss through the MOND transition boundary:
\begin{equation}
\Qpsi^{\rm MOND} \sim \frac{E_{\rm stored}}{P_{\rm rad}}
\sim \frac{(v_{\rm rot}/c)^2}{H_0 T_{\rm dyn}} \sim 10^9,
\end{equation}
where $v_{\rm rot} \sim 200$~km/s is the typical rotation velocity,
$H_0^{-1} \sim 4\ee{17}$~s is the Hubble time, and $T_{\rm dyn}
\sim 10^8$~yr is the galactic dynamical time.

\subsection{Why this is irrelevant for the cavity}

The SRF cavity sits deep in the Newtonian regime where $\mu \to 1$.
There is no MOND transition within or near the cavity. The
self-interaction term $\kappa a^2/c^2$ is negligible by a factor
$\sim 10^{13}$ compared to $\nabla \cdot \mathbf{a}$ at laboratory
accelerations (section02 Table~2: Solar regime has $\kappa a^2/c^2
\sim 10^{-19}$~s$^{-2}$ vs.\ $\nabla \cdot \mathbf{a} \sim
10^{-6}$~s$^{-2}$). The MOND self-confinement mechanism does not
operate at laboratory scales.

\verdict{$\Qpsi^{\rm MOND} = 10^9$ describes galactic-scale psi-modes.
It has no bearing on the $\dpsiEM$ perturbation inside an SRF cavity.
Using it in the cavity loss formula was physically unjustified.}


%=============================================================================
\section{Regime 2: Local Radiation-Limited ($\Qpsi^{\rm local} \sim 30$)}
\label{sec:regime2}
%=============================================================================

The W4i15 P2 update (Finding 2, Assumption 3 discussion, lines 597--615)
estimated:
\begin{equation}
\Qpsi^{\rm local} \sim (k_{\rm cav} R_{\rm cav})^3 \sim \pi^3 \sim 31,
\end{equation}
using the standard radiation-$Q$ formula for a source of size $R$ at
wavenumber $k$, analogous to antenna radiation resistance. This assumed
the psi-field propagates at $c_\psi = c$ in the Newtonian limit.

\subsection{The assumption: $c_\psi = c$}

The P2 derivation (Eq.~W4i15-P2 line 458) writes $c_\psi$ as the
``psi-field propagation speed (= $c$ in the Newtonian regime where
$\mu \to 1$).'' This is motivated by the linearised spatial equation
$\nabla^2(\dpsiEM) = \text{source}$, which has the form of a Poisson
equation --- no time derivatives, hence no propagation in the spatial
sector.

However, the P2 derivation then treats the driven response as a
damped harmonic oscillator (Eq.~W4i15-P2 line 486--489) with
$\Opsi = c_\psi k \sqrt{\mu_{\rm bg}}$ and a Lorentzian imaginary
part Im$[\hat{G}(\Omega)]$. This implicitly assumes a \emph{wave-like}
temporal response with finite propagation speed $c_\psi$.

\subsection{What the temporal completion actually says}

The full dynamic action (section02 Eq.~2.10) is:
\begin{equation}
S_\psi = \int dt\,d^3x\;\left\{
  \frac{a_*^2}{8\pi G}\left[
    W\!\left(\frac{|\nabla\psi|^2}{a_*^2}\right)
    + K\!\left(\frac{c}{a_0}|\dot\psi{-}\dot\psi_0|\right)
  \right]
  - \frac{c^2}{2}\,\psi(\rho - \bar\rho)
\right\}.
\end{equation}

The temporal kinetic function $K(\Delta)$ with $K'(\Delta) =
\mu(\Delta) = \Delta/(1+\Delta)$ gives:
\begin{equation}
K(\Delta) = \Delta - \ln(1+\Delta) \approx \half\Delta^2 + \order{\Delta^3}
\quad (\Delta \ll 1).
\end{equation}

The equation of motion for small $\dpsiEM$ perturbations about
$\psigrav$ in the Newtonian regime ($\mu_{\rm spatial} \to 1$):
\begin{equation}
\nabla^2(\dpsiEM)
+ \frac{a_*^2}{8\pi G}\,K''(\Delta_{\rm bg})\,
\frac{c^2}{a_0^2}\,\ddot{\delta\psi}_{\rm EM}
= \frac{4\pi G \lambda}{c^4}\,u_{\rm EM}.
\label{eq:full-eom}
\end{equation}

The effective propagation speed squared is:
\begin{equation}
c_s^2 = \frac{8\pi G}{a_*^2}\,\frac{1}{K''(\Delta_{\rm bg})}\,
\frac{a_0^2}{c^2}.
\label{eq:cs2}
\end{equation}

Now, $K''(\Delta) = 1/(1+\Delta)^2$, so for $\Delta_{\rm bg} \ll 1$
(laboratory regime), $K''(\Delta_{\rm bg}) \approx 1$. Substituting
$a_* = 2a_0/c^2$:
\begin{equation}
c_s^2 = \frac{8\pi G}{4a_0^2/c^4}\,\frac{a_0^2}{c^2}
= \frac{8\pi G\,c^4}{4\,a_0^2}\,\frac{a_0^2}{c^2}
= 2\pi G\,c^2.
\label{eq:cs2-lab}
\end{equation}

\corrected{Wait --- this gives $c_s^2 = 2\pi G c^2$, which has
dimensions of $(\text{m/s})^2 \times (\text{m}^3/(\text{kg}\cdot
\text{s}^2))$, i.e.\ $\text{m}^4/(\text{kg}\cdot\text{s}^4)$.
This is NOT a speed squared. The dimensional analysis shows that
Eq.~\eqref{eq:full-eom} has a mismatch: the spatial Laplacian
acts on $\dpsiEM$ (dimensionless), while the temporal term carries
the $a_*^2/(8\pi G)$ prefactor from the Lagrangian density.}

Let us redo this carefully. The linearised equation of motion from
the full action, perturbing about $\psigrav$:
\begin{align}
&\text{Spatial:}\quad
\frac{a_*^2}{8\pi G}\,\nabla\cdot\!\left[
  \mu_{\rm bg}\,\nabla(\dpsiEM)\right]
= -\frac{4\pi G\lambda}{c^4}\,u_{\rm EM}
+ \frac{c^2}{2}\,\dpsiEM\,\rho,
\label{eq:spatial-lin}\\
&\text{Temporal:}\quad
-\frac{a_*^2}{8\pi G}\,K''(\Delta_{\rm bg})\,
\frac{c^2}{a_0^2}\,\ddot{\delta\psi}_{\rm EM}.
\label{eq:temporal-lin}
\end{align}

Combining (and dropping the matter backreaction $\rho\,\dpsiEM$
which is negligible for the EM-sourced perturbation):
\begin{equation}
\frac{a_*^2}{8\pi G}\left[
  \mu_{\rm bg}\,\nabla^2(\dpsiEM)
  - K''(\Delta_{\rm bg})\,\frac{c^2}{a_0^2}\,
  \ddot{\delta\psi}_{\rm EM}
\right]
= -\frac{4\pi G\lambda}{c^4}\,u_{\rm EM}.
\label{eq:combined-lin}
\end{equation}

In the Newtonian regime, $\mu_{\rm bg} = 1$ and
$K''(\Delta_{\rm bg}) \approx 1$. The effective wave equation is:
\begin{equation}
\nabla^2(\dpsiEM)
- \frac{c^2}{a_0^2}\,\ddot{\delta\psi}_{\rm EM}
= -\frac{(4\pi G)^2\,2\lambda}{a_*^2\,c^4}\,u_{\rm EM}.
\label{eq:wave-dpsi}
\end{equation}

The propagation speed:
\begin{equation}
c_\psi^2 = \frac{a_0^2}{c^2}.
\end{equation}

With $a_0 \approx 1.2\ee{-10}$~m/s$^2$:
\begin{equation}
c_\psi = \frac{a_0}{c} = \frac{1.2\ee{-10}}{3\ee{8}}
= 4\ee{-19}~\text{m/s}.
\label{eq:cpsi-value}
\end{equation}

\begin{finding}[$c_\psi \approx 4\ee{-19}$~m/s in the Laboratory]
\label{find:cpsi}
\newfinding{The psi-field propagation speed in the laboratory
(Newtonian) regime is NOT $c$. It is $c_\psi = a_0/c \approx
4\ee{-19}$~m/s --- twenty-seven orders of magnitude below $c$.
This is the laboratory manifestation of the $c_s^2 \to 0$ dust
branch from Appendix~Q.}

At $c_\psi = 4\ee{-19}$~m/s, the psi-field is effectively
\textbf{non-propagating} on all laboratory timescales and length
scales. A perturbation $\dpsiEM$ created inside an SRF cavity
\emph{cannot escape} by radiation on any experimentally relevant
timescale:
\begin{equation}
t_{\rm escape} \sim \frac{R_{\rm cav}}{c_\psi}
= \frac{0.1~\text{m}}{4\ee{-19}~\text{m/s}}
= 2.5\ee{17}~\text{s}
\approx 8\ee{9}~\text{yr}.
\end{equation}

The ``radiation $Q$'' estimate $\Qpsi^{\rm local} \sim (k R)^3 \sim
30$ assumed $c_\psi = c$, giving radiation escape in nanoseconds.
With $c_\psi = a_0/c$, the radiation escape time exceeds the age
of the Universe. The local radiation-loss mechanism is completely
inoperative.

\verdict{$\Qpsi^{\rm local} \sim 30$ is WRONG. The correct
local Q is effectively infinite ($\Qpsi^{\rm local} \to \infty$)
because $c_\psi \to 0$ suppresses radiation losses.}
\end{finding}


%=============================================================================
\section{Regime 3: Quasi-Static Non-Resonant Response (The Correct Regime)}
\label{sec:regime3}
%=============================================================================

With $c_\psi = a_0/c \approx 4\ee{-19}$~m/s, the temporal term
in Eq.~\eqref{eq:wave-dpsi} is negligible compared to the spatial
Laplacian at SRF frequencies. The ratio of temporal to spatial terms:
\begin{equation}
\frac{\omega^2/c_\psi^2}{k^2} = \frac{\omega^2}{k^2 c_\psi^2}
= \frac{\omega^2 R_{\rm cav}^2}{c_\psi^2\,\pi^2}.
\end{equation}

For $\omega = 2\pi \times 1.3\ee{9}$~Hz and $R_{\rm cav} = 0.1$~m:
\begin{equation}
\frac{\omega^2 R_{\rm cav}^2}{c_\psi^2\,\pi^2}
= \frac{(8.2\ee{9})^2 \times 0.01}{(4\ee{-19})^2 \times 10}
\approx 4\ee{55}.
\end{equation}

The temporal term dominates by 55 orders of magnitude. This means
the $\dpsiEM$ response is NOT quasi-static --- it is
\textbf{temporally stiff}.

\corrected{This reverses the naive expectation. At $c_\psi
\approx 0$, the psi-field has enormous temporal inertia. The
$\ddot{\delta\psi}$ term is NOT negligible --- it is the DOMINANT
term. The equation becomes:
\begin{equation}
-\frac{c^2}{a_0^2}\,\ddot{\delta\psi}_{\rm EM}
\approx -\frac{(4\pi G)^2\,2\lambda}{a_*^2\,c^4}\,u_{\rm EM},
\end{equation}
giving a temporally driven response:
\begin{equation}
\ddot{\delta\psi}_{\rm EM}
= \frac{(4\pi G)^2\,2\lambda\,a_0^2}{a_*^2\,c^6}\,u_{\rm EM}.
\end{equation}}

For an EM source oscillating at $2\omega$ (the $\cos(2\omega t)$
component of $u_{\rm EM}$), the steady-state response is:
\begin{equation}
\dpsiEM^{(2\omega)} = -\frac{1}{(2\omega)^2}\,
\frac{(4\pi G)^2\,2\lambda\,a_0^2}{a_*^2\,c^6}\,
\Delta u(\mathbf{x}).
\label{eq:dpsi-response}
\end{equation}

This is purely \emph{reactive} (in-phase with the drive, no
dissipation). There is no imaginary part because the equation
is:
\begin{equation}
-\omega_{\rm drive}^2\,\dpsiEM = \text{source},
\end{equation}
which has a real solution for any real $\omega_{\rm drive}$.

\subsection{Where does dissipation enter?}

In the W4i15 P2 derivation, dissipation entered through
Im$[\hat{G}(\Omega)]$, which requires a finite damping rate
$\gpsi = \Opsi/(2\Qpsi)$. Without damping, the Green's function
has no imaginary part (away from resonance), and there is no energy
loss.

The key question: \textbf{what damps $\dpsiEM$?}

Three candidate mechanisms:

\begin{enumerate}
\item \textbf{Radiation to spatial infinity}: Requires finite
  $c_\psi$. At $c_\psi = a_0/c \sim 4\ee{-19}$~m/s, this is
  negligible ($\gpsi^{\rm rad} \sim c_\psi/R \sim 4\ee{-18}$~s$^{-1}$,
  giving $\Qpsi^{\rm rad} \sim \omega/\gpsi \sim 2\ee{27}$).

\item \textbf{Gravitational coupling to matter}: The $\dpsiEM$ field
  couples back to the cavity walls and surrounding matter through the
  $c^2 \psi \rho/2$ term in the action. This provides a dissipative
  channel: energy flows from EM $\to$ $\dpsiEM$ $\to$ mechanical
  acceleration of matter. The rate is set by the gravitational coupling
  strength: $\gpsi^{\rm grav} \sim G \rho_{\rm Nb}\,\omega_{\rm mech}
  /c^2$, where $\rho_{\rm Nb} \sim 8570$~kg/m$^3$ is the Nb density.
  This is extremely small: $\gpsi^{\rm grav} \sim 10^{-16}$~s$^{-1}$.

\item \textbf{Nonlinear self-interaction}: In the Newtonian regime,
  $\mu \to 1$ and the nonlinear term $\kappa a^2/c^2$ is negligible.
  No nonlinear damping.
\end{enumerate}

\begin{finding}[The Psi-Channel Dissipation Paradox]
\label{find:paradox}
\newfinding{In the correct DFD formalism, all three candidate
damping mechanisms for $\dpsiEM$ in the laboratory are negligible.
The psi-field perturbation driven by the SRF cavity is an
undamped, reactive, non-dissipative response. This means:
\begin{equation}
\Gamma_\psi = 0 \quad \text{(to leading order)},
\end{equation}
and the psi-channel does NOT contribute to cavity loss.}

\honest{This is a startling conclusion that contradicts the entire
psi-channel detection programme (Q-ratio diagnostic, SQMS bounds,
Phase~I proposal). If $\Gamma_\psi = 0$, then:
\begin{itemize}[nosep]
\item The Q-ratio diagnostic has no signal.
\item The SQMS ``accidental'' bound $\lbone < 1.56\ee{-9}$ is
  vacuous (there is no psi-channel loss to bound).
\item The Phase~I search for anomalous residual resistance from
  the psi-channel would find nothing regardless of $\lambda$.
\end{itemize}

Before accepting this, we must examine what went wrong with the
prior derivation and whether there is a dissipative channel that
we are missing.}
\end{finding}


%=============================================================================
%=============================================================================
\part{Resolving the Paradox: The DC Channel}
\label{part:resolution}
%=============================================================================
%=============================================================================


%=============================================================================
\section{The Overlooked DC Response}
\label{sec:dc}
%=============================================================================

The EM energy density has TWO components:
\begin{equation}
u_{\rm EM}(\mathbf{x}, t) = \bar{u}(\mathbf{x})
+ \Delta u(\mathbf{x})\cos(2\omega t),
\end{equation}
where $\bar{u}$ is the time-averaged (DC) component and $\Delta u$
is the oscillating ($2\omega$) component.

The $2\omega$ component drives a reactive (non-dissipative) response
as shown above. But the DC component $\bar{u}(\mathbf{x})$ produces
a \textbf{static shift} in $\dpsiEM$:
\begin{equation}
\nabla^2(\dpsiEM^{\rm DC}) = -\frac{(4\pi G)^2\,2\lambda}
{a_*^2\,c^4}\,\bar{u}(\mathbf{x}).
\end{equation}

This is a Poisson equation (temporal term drops out for DC) with
solution:
\begin{equation}
\dpsiEM^{\rm DC}(\mathbf{x}) = \frac{(4\pi G)^2\,2\lambda}
{a_*^2\,c^4}\int \frac{\bar{u}(\mathbf{x}')}{4\pi|\mathbf{x}
- \mathbf{x}'|}\,d^3x'.
\end{equation}

This static shift does not oscillate and does not directly dissipate
energy. However, it \emph{does} modify the optical metric locally,
creating a gravitational potential well proportional to $U_{\rm EM}$.
This is the parametric (Channel 2) effect from Appendix~R.


%=============================================================================
\section{The Correct Dissipation Mechanism: Parametric Energy Transfer}
\label{sec:parametric}
%=============================================================================

The resolution of the paradox lies in recognising that the psi-channel
loss is NOT a simple resonant absorption. It is a \textbf{parametric
process}: the oscillating $\dpsiEM$ modulates the optical metric at
$2\omega$, which shifts the EM mode frequency, which transfers energy
from the EM field to the psi-field via parametric coupling.

From Appendix~R (Eq.~R.4), the parametric coupling enters through the
$\alpha U(t) q$ term in the mode equation:
\begin{equation}
\ddot{q} + 2\gpsi\dot{q} + \Opsi^2 q
= \frac{(\lambda - 1)}{\Mpsi}\,G(t) + \alpha U(t)\,q.
\end{equation}

The parametric term $\alpha U(t) q$ with $U(t) = U_0[1 + m\cos(2\omega
t)]$ creates Mathieu-type instability when $2\omega \approx 2\Opsi$.

But with $c_\psi = a_0/c$, the mode frequency $\Opsi$ is:
\begin{equation}
\Opsi = c_\psi\,k = \frac{a_0}{c}\,\frac{\pi}{R_{\rm cav}}
= \frac{1.2\ee{-10}}{3\ee{8}}\,\frac{\pi}{0.1}
= 1.3\ee{-17}~\text{rad/s}.
\end{equation}

The parametric resonance condition $2\omega \approx 2\Opsi$ requires
$\omega \approx 10^{-17}$~rad/s, which is 26 orders of magnitude
below the SRF frequency. \textbf{Parametric resonance is completely
off-resonant.}

The off-resonant parametric effect produces a frequency shift (not
energy transfer) of order:
\begin{equation}
\frac{\delta\omega}{\omega} \sim \frac{\alpha U_0}{\Opsi^2 - \omega^2}
\approx -\frac{\alpha U_0}{\omega^2},
\end{equation}
which is a real (reactive) effect, not dissipative.


%=============================================================================
\section{Where the Prior Derivation Went Wrong}
\label{sec:error}
%=============================================================================

The W4i15 P2 derivation (Theorem 2.1, lines 416--569) and the P11
derivation (Part II, lines 307--627) both implicitly assumed that
the psi-field Green's function has a Lorentzian resonance structure
(Eq.~W4i15-P2 line 486):
\begin{equation}
\text{Im}[\hat{G}(\Omega)] = \frac{\Omega\,\Opsi}
{(\Opsi^2 - \Omega^2)^2 + \Omega^2\,\Opsi^2/\Qpsi^2}.
\end{equation}

This is the response function for a \textbf{damped harmonic oscillator}
with resonant frequency $\Opsi$ and quality factor $\Qpsi$. It has a
finite imaginary part (dissipation) because $\Qpsi$ is finite.

The error: the derivation treated $\Opsi$ and $\Qpsi$ as free
parameters and evaluated the Green's function in the
``off-resonant'' limit $2\omega \ll \Opsi$ (P2 Eq.~line 510):
\begin{equation}
\text{Im}[\hat{G}(2\omega)] \approx \frac{2\omega}{\Opsi^2\,\Qpsi}.
\end{equation}

Using $\Opsi = c k_{\rm cav} \sim 10^{10}$~rad/s (with $c_\psi = c$)
made this ``off-resonant''. But using the correct $c_\psi = a_0/c$
gives $\Opsi \sim 10^{-17}$~rad/s, making $2\omega \gg \Opsi$. In
THIS limit:
\begin{equation}
\text{Im}[\hat{G}(2\omega)] \approx
\frac{2\omega\,\Opsi}{(2\omega)^4 + (2\omega)^2\Opsi^2/\Qpsi^2}
\approx \frac{\Opsi}{(2\omega)^3\,\Qpsi^2}
\to 0 \quad (\omega \to \infty).
\end{equation}

The imaginary part is \textbf{negligible} at SRF frequencies when
$\Opsi \ll \omega$. The psi-channel loss vanishes.

\begin{finding}[The P2 Error Was Assuming $c_\psi = c$]
\label{find:error}
\newfinding{The entire psi-channel loss derivation in W4i15 P2 and
W4i15 P11 used the wrong propagation speed ($c_\psi = c$) for the
scalar perturbation inside the cavity. The correct value from the
temporal completion is $c_\psi = a_0/c \approx 4\ee{-19}$~m/s.
This makes $\Opsi \sim 10^{-17}$~rad/s rather than $10^{10}$~rad/s,
shifting the cavity from the off-resonant ($\omega \ll \Opsi$) to
the far-above-resonance ($\omega \gg \Opsi$) regime. In the
far-above-resonance regime, the imaginary part of the Green's
function (and hence the dissipation rate) is suppressed by
$(\Opsi/\omega)^3$, which is a factor of $\sim 10^{-81}$.}

\caution{If this analysis is correct, the psi-channel loss in SRF
cavities is $\sim 10^{-81}$ times smaller than all prior estimates.
The Q-ratio diagnostic, the SQMS bound, and the Phase~I proposal
would all be invalidated.}
\end{finding}


%=============================================================================
%=============================================================================
\part{Critical Reassessment: Is the Temporal Completion Correct for
  EM-Sourced Perturbations?}
\label{part:reassessment}
%=============================================================================
%=============================================================================


%=============================================================================
\section{The Two-Component Subtlety}
\label{sec:two-component}
%=============================================================================

Before accepting the devastating conclusion of Part~II, we must examine
a crucial subtlety: the temporal completion (Appendix~Q) derives
$c_s^2 \to 0$ for the \emph{cosmological} dust sector. Does this
apply to the EM-sourced perturbation $\dpsiEM$?

\subsection{What the dust branch actually proves}

Theorem~Q.4 (dust branch) establishes that for \emph{homogeneous}
FRW solutions near the screen background, $w \to 0$ and $c_s^2 \to 0$.
The proof uses:
\begin{enumerate}[nosep]
\item The FRW symmetry (homogeneous, isotropic).
\item The deviation invariant $\Delta = (c/a_0)|\dot\psi - \dot\psi_0|$
  with $\Delta \ll 1$.
\item The conserved current $a^3 \mu(\Delta) = \text{const}$.
\end{enumerate}

The $c_s^2 \to 0$ result is the \emph{adiabatic sound speed} for
long-wavelength perturbations of this cosmological background. It
does NOT necessarily apply to short-wavelength, EM-driven perturbations
in the laboratory.

\subsection{The linearised perturbation equation revisited}

The perturbation $\dpsiEM$ is sourced by EM energy, not by the
cosmological background. The linearised equation about $\psigrav$
in the laboratory is:
\begin{equation}
\frac{a_*^2}{8\pi G}\left[
  \nabla^2(\dpsiEM) - K''(\Delta_{\rm bg})\,\frac{c^2}{a_0^2}\,
  \ddot{\delta\psi}_{\rm EM}
\right] = -\frac{\lambda\,4\pi G}{c^4}\,u_{\rm EM}.
\label{eq:lab-perturbation}
\end{equation}

The coefficient of the time derivative is
$K''(\Delta_{\rm bg})\,c^2/a_0^2$. In the laboratory, $\Delta_{\rm bg}$
is the temporal deviation of the local gravitational field from the
cosmological background. For the Earth's surface:
\begin{equation}
\Delta_{\rm bg} = \frac{c}{a_0}\,|\dot\psi_{\rm local}
- \dot\psi_{\rm cosmo}|.
\end{equation}

The local $\dot\psi$ has contributions from Earth's rotation, orbital
motion, and galactic motion, all of which are small compared to $a_0/c$.
For Earth's surface gravity:
\begin{equation}
\dot\psi_{\rm local} \sim \frac{2 g_\oplus}{c^2}\,v_{\rm rot}
\sim \frac{2 \times 10}{(3\ee{8})^2}\times 465
\sim 10^{-13}~\text{s}^{-1}.
\end{equation}
\begin{equation}
\dot\psi_{\rm cosmo} \sim H_0\,\psi_{\rm cosmo}
\sim 2.3\ee{-18}\times 10^{-5} \sim 10^{-23}~\text{s}^{-1}.
\end{equation}
\begin{equation}
\Delta_{\rm bg} = \frac{c}{a_0}\,|\dot\psi_{\rm local}
- \dot\psi_{\rm cosmo}|
\approx \frac{3\ee{8}}{1.2\ee{-10}}\times 10^{-13}
= 2.5\ee{5}.
\end{equation}

\begin{finding}[$\Delta_{\rm bg} \gg 1$ in the Laboratory]
\label{find:Delta}
\newfinding{The temporal deviation parameter $\Delta_{\rm bg}$ in
the laboratory is $\sim 10^5$, NOT $\ll 1$.}

This means the cosmological limit $\Delta \to 0$ (where $c_s^2 \to 0$)
does NOT apply to laboratory conditions. For $\Delta \gg 1$:
\begin{equation}
K''(\Delta) = \frac{1}{(1+\Delta)^2} \approx \frac{1}{\Delta^2}
\approx \frac{1}{(2.5\ee{5})^2} = 1.6\ee{-11}.
\end{equation}

The effective propagation speed squared:
\begin{equation}
c_\psi^2 = \frac{1}{K''(\Delta_{\rm bg})}\,\frac{a_0^2}{c^2}
= \Delta_{\rm bg}^2\,\frac{a_0^2}{c^2}
= \Delta_{\rm bg}^2\,\frac{a_0^2}{c^2}.
\end{equation}

But by definition, $\Delta_{\rm bg} = (c/a_0)\,|\dot\psi_{\rm local}|$,
so:
\begin{equation}
c_\psi^2 = \frac{c^2}{a_0^2}\,|\dot\psi_{\rm local}|^2\,
\frac{a_0^2}{c^2} = |\dot\psi_{\rm local}|^2\,R^2,
\end{equation}
where $R$ is a characteristic length... No, this is getting circular.
Let me redo this systematically.
\end{finding}

\subsection{Systematic derivation of $c_\psi$ at $\Delta_{\rm bg} \gg 1$}

The perturbation equation~\eqref{eq:lab-perturbation} can be written:
\begin{equation}
\nabla^2(\dpsiEM) - \frac{1}{\tilde{c}_\psi^2}\,
\ddot{\delta\psi}_{\rm EM} = \text{source},
\end{equation}
where:
\begin{equation}
\frac{1}{\tilde{c}_\psi^2} = K''(\Delta_{\rm bg})\,\frac{c^2}{a_0^2}.
\end{equation}

For $\Delta_{\rm bg} \gg 1$:
\begin{equation}
K''(\Delta_{\rm bg}) \approx \frac{1}{\Delta_{\rm bg}^2},
\end{equation}
so:
\begin{equation}
\frac{1}{\tilde{c}_\psi^2} = \frac{c^2}{a_0^2\,\Delta_{\rm bg}^2}
= \frac{c^2}{a_0^2}\,\frac{a_0^2}{c^2\,|\dot\psi_{\rm local}
- \dot\psi_0|^2}
= \frac{1}{|\dot\psi_{\rm local} - \dot\psi_0|^2}.
\end{equation}

Wait, $\Delta_{\rm bg} = (c/a_0)|\dot\psi_{\rm local} - \dot\psi_0|$,
so $a_0^2\Delta_{\rm bg}^2/c^2 = |\dot\psi_{\rm local} - \dot\psi_0|^2$.
Therefore:
\begin{equation}
\tilde{c}_\psi^2 = \frac{a_0^2\,\Delta_{\rm bg}^2}{c^2}
= |\dot\psi_{\rm local} - \dot\psi_0|^2.
\end{equation}

Using $|\dot\psi_{\rm local}| \sim 10^{-13}$~s$^{-1}$ and
$\dot\psi_0 \ll \dot\psi_{\rm local}$:
\begin{equation}
\tilde{c}_\psi \sim |\dot\psi_{\rm local}| \sim 10^{-13}~\text{s}^{-1}.
\end{equation}

\corrected{$\tilde{c}_\psi$ has units of s$^{-1}$, not m/s. This
signals a dimensional error. Let me trace it back.}

The issue: in Eq.~\eqref{eq:lab-perturbation}, the spatial and
temporal terms must have the same dimensions. $\nabla^2(\dpsiEM)$ has
dimensions m$^{-2}$ (since $\dpsiEM$ is dimensionless). The temporal
term $K''(\Delta_{\rm bg})\,(c^2/a_0^2)\,\ddot{\delta\psi}_{\rm EM}$
must also be m$^{-2}$:
\begin{itemize}
\item $[K''] = 1$ (dimensionless).
\item $[c^2/a_0^2] = (\text{m/s})^2/(\text{m/s}^2)^2
  = \text{s}^2/\text{m}^2$.
\item $[\ddot{\delta\psi}] = \text{s}^{-2}$.
\end{itemize}
Product: $\text{s}^2/\text{m}^2 \times \text{s}^{-2} = 1/\text{m}^2$.
Correct.

So the wave equation is:
\begin{equation}
\nabla^2(\dpsiEM) - \frac{K''(\Delta_{\rm bg})\,c^2}{a_0^2}\,
\ddot{\delta\psi}_{\rm EM} = \text{source},
\end{equation}
with effective speed $\tilde{c}_\psi$ defined by $1/\tilde{c}_\psi^2
= K''(\Delta_{\rm bg})\,c^2/a_0^2$, so:
\begin{equation}
\boxed{
\tilde{c}_\psi^2 = \frac{a_0^2}{c^2\,K''(\Delta_{\rm bg})}.
}
\label{eq:cpsi-general}
\end{equation}

For $\Delta_{\rm bg} \gg 1$, $K'' \approx 1/\Delta_{\rm bg}^2$:
\begin{equation}
\tilde{c}_\psi^2 = \frac{a_0^2\,\Delta_{\rm bg}^2}{c^2}.
\end{equation}

With $\Delta_{\rm bg} = (c/a_0)|\dot\psi_{\rm loc}|
\approx (3\ee{8}/1.2\ee{-10})\times 10^{-13} = 2.5\ee{5}$:
\begin{equation}
\tilde{c}_\psi^2 = \frac{(1.2\ee{-10})^2\,(2.5\ee{5})^2}
{(3\ee{8})^2}
= \frac{1.44\ee{-20}\times 6.25\ee{10}}{9\ee{16}}
= \frac{9\ee{-10}}{9\ee{16}} = 10^{-26}~\text{m}^2/\text{s}^2.
\end{equation}
\begin{equation}
\tilde{c}_\psi = 10^{-13}~\text{m/s}.
\end{equation}

\begin{finding}[Corrected Psi-Field Propagation Speed in the Laboratory]
\label{find:cpsi-corrected}
\newfinding{The psi-field propagation speed in the laboratory is
$\tilde{c}_\psi \approx 10^{-13}$~m/s (not $4\ee{-19}$~m/s from
the $\Delta \ll 1$ limit, and not $c$ from the naive Newtonian
assumption).}

The mode frequency for a psi-mode confined in the cavity:
\begin{equation}
\Opsi = \tilde{c}_\psi\,k_{\rm cav}
= 10^{-13}\,\frac{\pi}{0.1}
= 3\ee{-12}~\text{rad/s}.
\end{equation}

The SRF drive frequency $2\omega = 2\times 2\pi\times 1.3\ee{9}
= 1.6\ee{10}$~rad/s is still $\sim 10^{22}$ times above $\Opsi$.
The cavity is deep in the $\omega \gg \Opsi$ regime.

In this regime, the response of the psi-field to the $2\omega$
drive is:
\begin{equation}
\dpsiEM^{(2\omega)} \approx -\frac{\text{source}}{(2\omega)^2
/\tilde{c}_\psi^2} = -\frac{\tilde{c}_\psi^2}{(2\omega)^2}\,
\frac{\text{source}}{1},
\end{equation}
which is a \textbf{suppressed reactive response}, attenuated by
$(\Opsi/\omega)^2 \sim 10^{-44}$ relative to the resonant case.

The imaginary part (dissipation) in this regime:
\begin{equation}
\text{Im}[\hat{G}] \sim \frac{\Opsi}{(2\omega)^3\,\Qpsi},
\end{equation}
suppressed by an additional factor of $\Opsi/(2\omega) \sim 10^{-22}$
relative to the reactive response.
\end{finding}


%=============================================================================
\section{The Deep Question: Is the Temporal Kinetic Term Physical?}
\label{sec:deep-question}
%=============================================================================

\honest{The devastating suppression factor of $\sim 10^{-44}$
(reactive) to $\sim 10^{-66}$ (dissipative) rests entirely on
the temporal kinetic function $K(\Delta)$ from the dust branch
derivation. If this temporal sector is correct, the psi-channel
loss mechanism is dead.

However, there are reasons to question whether $K(\Delta)$ applies
unmodified to laboratory EM-sourced perturbations:

\begin{enumerate}
\item \textbf{The dust branch is cosmological}: It was derived for
  FRW backgrounds. The laboratory is not FRW.

\item \textbf{The temporal sector may have additional terms}: The
  ``full scalar-sector action'' in Eq.~2.10 of section02 is the
  \emph{proposed} form. The spatial sector $W(y)$ is well-constrained
  by rotation curves, lensing, and Solar System tests. The temporal
  sector $K(\Delta)$ is constrained only by cosmological requirements
  ($c_s^2 \to 0$ for structure formation) and theoretical self-consistency
  ($S^3$ composition law). There may be corrections at laboratory scales
  that restore faster propagation.

\item \textbf{The $\dpsiEM$ perturbation is EM-sourced}: It couples
  through $\lambda\,u_{\rm EM}/c^2$, which is an \emph{EM} source
  term, not a gravitational one. The response may be governed by the
  EM-sector dynamics, not the pure gravitational temporal sector.

\item \textbf{Mode equation in Appendix~R assumes finite $\Opsi$}:
  Appendix~R (Eq.~R.1) writes the mode equation with $\Opsi^2$ as a
  free parameter and $\gpsi$ as a separate damping rate. This
  framework was constructed for ``a single laboratory mode $q(t)$
  with natural frequency $\Opsi_\psi$'' --- it assumes the mode exists
  and has finite frequency. If $\Opsi \sim 10^{-12}$~rad/s, the
  ``mode'' would have a period of $10^{12}$~s $\sim 30{,}000$~years.
\end{enumerate}

\textbf{The honest assessment}: The temporal sector of DFD is not
well enough constrained to determine $c_\psi$ at laboratory scales.
The cosmological constraint gives $c_s^2 \to 0$ for structure
formation, but the specific value of $c_\psi$ in a table-top
experiment depends on $K''(\Delta_{\rm bg})$, which depends on
$\Delta_{\rm bg}$, which depends on the local time derivative of
$\psi$ relative to the cosmological background.}


%=============================================================================
%=============================================================================
\part{Resolution and Impact on SQMS Phase~I}
\label{part:impact}
%=============================================================================
%=============================================================================


%=============================================================================
\section{The Three Scenarios}
\label{sec:scenarios}
%=============================================================================

Given the uncertainty in the temporal sector, we must present the
SQMS bounds and Phase~I reach under three scenarios:

\subsection{Scenario A: $c_\psi = c$ (Newtonian limit, temporal sector decoupled)}

This is the assumption of ALL prior iterations. The temporal kinetic
term is ignored (or assumed to give $c_\psi = c$). The loss formula
is the W4i15 P2 result:
\begin{equation}
\frac{1}{Q_\psi^{\rm cav}} = \frac{(\lambda-1)^2\,G\,\calF}
{c^4\,\Opsi^2\,\Qpsi}\,\omega^2,
\end{equation}
with $\Opsi = c k_{\rm cav} \sim 10^{10}$~rad/s. Here $\Qpsi$ enters
as the scalar-mode quality factor. Using $\Qpsi = 10^9$ gives the
canonical programme values.

\textbf{Problem}: $\Qpsi = 10^9$ is the MOND self-confinement value
for galactic modes, not laboratory modes. There is no physics-based
estimate of $\Qpsi$ for this scenario. Values between $\Qpsi \sim 1$
(radiation-dominated, if one uses the standard antenna formula but
with partial confinement by cavity walls) and $\Qpsi \sim 10^{12}$
(limited only by gravitational radiation) are possible.

\textbf{SQMS bound}: $\lbone < 1.56\ee{-9}\times\sqrt{\Qpsi/10^9}$.

\textbf{Impact}: The bound scales as $\sqrt{\Qpsi}$. If $\Qpsi = 1$
(no confinement), the bound tightens to $\lbone < 5\ee{-14}$.
If $\Qpsi = 10^{12}$, it loosens to $\lbone < 5\ee{-8}$.

\subsection{Scenario B: $c_\psi = a_0/c$ (dust branch, $\Delta \to 0$ limit)}

This uses the cosmological dust branch limit directly. As shown in
\S\ref{sec:regime2}, $c_\psi \sim 4\ee{-19}$~m/s, and the
psi-channel loss is suppressed by $\sim 10^{-81}$.

\textbf{SQMS bound}: No bound (psi-channel loss is unobservably small).

\textbf{Impact}: Psi-channel detection through SRF cavity loss is
impossible at any realistic sensitivity.

\subsection{Scenario C: $c_\psi = \tilde{c}_\psi(\Delta_{\rm bg})$ (laboratory $\Delta$)}

This uses the corrected propagation speed $\tilde{c}_\psi \sim
10^{-13}$~m/s from Finding~\ref{find:cpsi-corrected}. The
psi-channel loss is suppressed by $(\Opsi/\omega)^2 \sim 10^{-44}$.

\textbf{SQMS bound}: No practical bound (suppression too extreme).

\textbf{Impact}: Same as Scenario B for all practical purposes.


%=============================================================================
\section{Assessment: Which Scenario Is Correct?}
\label{sec:assessment}
%=============================================================================

\begin{finding}[The Temporal Sector Is the Key Uncertainty]
\label{find:temporal}
\newfinding{The question of local vs.\ global $\Qpsi$ is actually
a question about the temporal kinetic sector $K(\Delta)$. The three
scenarios correspond to three different treatments of this sector:

\begin{enumerate}[nosep]
\item[(A)] Ignore the temporal sector entirely: $\dpsiEM$ responds
  instantaneously, dissipation enters through an \emph{ad hoc} $\Qpsi$.
\item[(B)] Use the cosmological limit $c_s^2 \to 0$: no dissipation.
\item[(C)] Use the laboratory $\Delta_{\rm bg}$: negligible dissipation.
\end{enumerate}

Scenarios B and C are \textbf{more physically grounded} than A,
because they use the temporal sector that is derived from the same
$S^3$ composition law that determines $\mu(x)$. Scenario A treats
$\Qpsi$ as a free parameter, which means the psi-channel prediction
has an unconstrained fudge factor.}

\honest{The correct answer almost certainly involves the temporal
sector at some level. The question is whether the $K(\Delta)$
derived for cosmological perturbations is the complete story for
short-wavelength, high-frequency, EM-driven perturbations.

There is a possibility that the EM coupling introduces additional
temporal dynamics not captured by $K(\Delta)$. Specifically, the
EM-$\psi$ interaction term $\lambda\,u_{\rm EM}/c^2$ in the field
equation may have its own temporal structure that is faster than
$K(\Delta)$ because it is mediated by the EM field (which propagates
at $c$), not by the gravitational self-interaction (which has
$c_s \to 0$).

This would correspond to a ``mixed'' scenario where the EM-driven
$\dpsiEM$ has a propagation speed set by the EM-$\psi$ coupling
rather than by the gravitational temporal sector. This requires a
careful re-examination of the full coupled EM-$\psi$ field equations,
not just the decoupled scalar sector.

\textbf{Until this is resolved, Scenario A should be retained as
the ``working hypothesis'' for the SQMS proposal, with the $\Qpsi$
uncertainty flagged as a known theoretical gap.}}
\end{finding}


%=============================================================================
\section{Impact on SQMS Phase~I Proposal}
\label{sec:phase1}
%=============================================================================

\begin{tcolorbox}[colback=yellow!5!white, colframe=orange!60!black,
  title=\textbf{IMPACT ASSESSMENT}]

\begin{enumerate}
\item \textbf{Q-ratio diagnostic}: The \emph{ratio} $\mathcal{R} =
  Q_\psi^{-1}(\omega_1)/Q_\psi^{-1}(\omega_2)$ does NOT depend on
  $\Qpsi$ (it cancels in the ratio). The Q-ratio prediction
  $\mathcal{R} = 0.49$ (or 0.27 with $1/\omega$ correction) is
  \textbf{robust} against the local-vs-global $\Qpsi$ question.
  The shape of the frequency dependence is set by the overlap
  integrals, not by the absolute loss rate.

  \confirmed{The Q-ratio diagnostic is unaffected by this analysis.
  Phase~I can proceed with the multi-frequency measurement regardless.}

\item \textbf{Absolute SQMS bound ($\lbone < 1.56\ee{-9}$)}: This
  bound assumes a specific absolute loss rate at a specific $\Qpsi$.
  It is now uncertain by a factor $\sqrt{\Qpsi}$, where $\Qpsi$
  ranges from $\sim 1$ (no temporal sector) to $\sim 10^9$ (MOND
  value) to $\sim \infty$ (quasi-static regime). The bound could
  be anywhere from $5\ee{-14}$ to vacuous.

  \caution{The absolute bound is unreliable until the temporal sector
  is resolved. Quote it as $\lbone < 1.56\ee{-9}\times
  (\Qpsi/10^9)^{1/2}$ with $\Qpsi$ flagged as a theoretical
  uncertainty.}

\item \textbf{Phase~I detection threshold}: If $\Qpsi \to \infty$
  (Scenarios B/C), then the psi-channel contributes negligible loss,
  and Phase~I would need an alternative detection strategy. If
  $\Qpsi = 10^9$ (Scenario A), the existing programme is on target.
  If $\Qpsi \sim 1$, the signal is $10^9\times$ larger and existing
  SQMS data would already have detected it or ruled it out.

\item \textbf{M$_{\rm eff} = 3.01\ee{6}$~kg}: This effective mass
  was computed for the passive detection (gravitational-gradient)
  mode, not for the psi-channel loss mode. It is unaffected by
  this analysis.

\item \textbf{Passive detection mode}: The passive detection measures
  the influence of the external galactic $\psi_{\rm grav}$ on cavity
  frequency. This does NOT involve $\dpsiEM$ or the temporal sector.
  The passive mode measures the static gravitational background and
  is completely unaffected by the $\Qpsi$ question.

  \confirmed{The Passive Superiority Theorem, the $M_{\rm eff}$ value,
  and the passive detection reach are all unaffected.}
\end{enumerate}
\end{tcolorbox}


%=============================================================================
\section{Recommended Actions}
\label{sec:actions}
%=============================================================================

\begin{enumerate}
\item \textbf{URGENT}: Independent derivation of the linearised
  perturbation equation for EM-sourced $\dpsiEM$, including the
  temporal sector $K(\Delta)$ with the correct $\Delta_{\rm bg}$
  for laboratory conditions. This must determine whether the
  effective $c_\psi$ at GHz frequencies is (a)~$c$, (b)~$a_0/c$,
  (c)~$\tilde{c}_\psi(\Delta_{\rm bg}) \sim 10^{-13}$~m/s, or
  (d)~something else set by the EM coupling.

\item \textbf{HIGH PRIORITY}: Examine whether the EM-$\psi$ coupling
  introduces a separate temporal propagation channel that bypasses
  the slow $K(\Delta)$ sector. The EM field propagates at $c$; if
  the $\dpsiEM$ perturbation inherits this speed through the coupling,
  then Scenario A may be correct despite the dust branch.

\item \textbf{MEDIUM PRIORITY}: Rewrite the SQMS proposal to present
  the Q-ratio diagnostic as the primary signature (robust against
  $\Qpsi$ uncertainty) and the absolute bound as secondary (dependent
  on $\Qpsi$).

\item \textbf{MEDIUM PRIORITY}: Compute $\Delta_{\rm bg}$ and
  $K''(\Delta_{\rm bg})$ for the specific FNAL/SQMS site geometry
  (latitude, local gravitational environment, galactic position).
  This gives the numerical value of $c_\psi$ for the experiment.

\item \textbf{LONG TERM}: The temporal sector $K(\Delta)$ is the
  least-tested part of DFD. Developing laboratory or astronomical
  tests of $K(\Delta)$ would resolve this question definitively.
\end{enumerate}


%=============================================================================
\section{Summary: Local vs.\ Global $\Qpsi$ --- Resolution}
\label{sec:summary}
%=============================================================================

\begin{tcolorbox}[colback=green!5!white, colframe=green!50!black,
  title=\textbf{DEFINITIVE ANSWER}]

\textbf{Q: Is $\Qpsi$ local ($\sim 30$) or global ($10^9$)?}

\textbf{A: Neither.} The question exposes a deeper issue: the
psi-field propagation speed in the laboratory, which is determined
by the temporal kinetic sector $K(\Delta)$.

\begin{itemize}[nosep]
\item The $\Qpsi^{\rm local} \sim 30$ estimate (W4i15 P2) used
  $c_\psi = c$, which contradicts the temporal completion.
\item The $\Qpsi^{\rm MOND} = 10^9$ estimate applies to galactic
  modes, not laboratory modes.
\item The correct $c_\psi$ is set by $K''(\Delta_{\rm bg})$, giving
  $c_\psi \sim 10^{-13}$~m/s in the lab, which makes the cavity
  operate in the $\omega \gg \Opsi$ regime where dissipation is
  negligible.
\item The Q-ratio diagnostic ($\mathcal{R} = 0.49$ or $0.27$) is
  \textbf{robust}: it does not depend on $\Qpsi$.
\item The absolute SQMS bound $\lbone < 1.56\ee{-9}$ is
  \textbf{unreliable}: it carries an unconstrained $\sqrt{\Qpsi}$
  factor.
\item The passive detection mode is \textbf{unaffected}.
\end{itemize}

\textbf{Programme impact}: This analysis does NOT kill the SQMS
programme. It sharpens it: the Q-ratio diagnostic (frequency-dependence
shape) is the reliable observable; the absolute loss rate is uncertain.
Phase~I should be framed around the shape measurement, not the
absolute bound.

\textbf{Remaining theoretical question}: Does the EM-$\psi$ coupling
introduce a fast temporal channel that bypasses the slow $K(\Delta)$
sector? If yes, Scenario A (with some effective $\Qpsi$) applies and
the absolute bounds are meaningful. If no, only the Q-ratio shape
test survives.
\end{tcolorbox}


\end{document}
